\section{The remaining clearing exponents}\label{app:small}
The isolated-pole estimate handles the large odd primes. Here we
supply the bounds at $2$ and at odd primes with $p^2\le7n$.
These estimates follow the local-series and Legendre bounds in
\cite[Sections 4.1 and 4.5]{krisman}. Their purpose is to show that
the small odd primes cost only $o(n)$ after factorial normalization,
whereas the prime $2$ contributes the term $15\log2$.

\subsection{Clearing the coefficients}
Put $Z=3\#\Lambda_n=\tfrac{15}{2}n+6$. For odd $p$, define
\[
 \lambda_p=\lfloor\log_p(7n)\rfloor,\qquad
 \nu_p(a)=3\sum_{\ell\in\Lambda_n}v_p(2a+1+2\ell)
                         -\sum_{b\ne a}B_bv_p(b-a).
\]
At $2$, use
\[
 \lambda_2=\lfloor\log_2n\rfloor,\qquad
 \nu_2(a)=1-Z-\sum_{b\ne a}B_bv_2(b-a).
\]
The factor $n-2a+2u$ in the local series \eqref{eq:local-series}
is integral at odd primes, so it can be omitted in a lower bound.
At $2$, its constant and linear terms are divisible by $2$.
Each half-integer numerator constant has valuation $-1$, giving
the term $-Z$ in $\nu_2$.

Every nonzero numerator integer or pole difference has absolute value
below $7n$, and each pole difference has absolute value at most $n$.
The series rule from Lemma~\ref{lem:local}, now with slope
$\lambda_p$, therefore gives
\begin{equation}\label{eq:small-local}
 v_p(a_{i,a})\ge\nu_p(a)-\lambda_p(B_a-i).
\end{equation}
The harmonic sums satisfy
$v_p(\sigma_a(i+2))\ge-(i+2)\lambda_p$ at odd $p$ and are
integral at $2$. Hence the integer exponents
\begin{align}
 e_p&=\max\left(0,\max_{0\le a\le n}(-\nu_p(a)+33\lambda_p)\right)
                &&(p\text{ odd},\ p^2\le7n),                 \label{eq:small-ep}\\
 e_2&=\max\left(0,\max_{0\le a\le n}(-\nu_2(a)+31\lambda_2)\right)
                                                               \label{eq:two-ep}
\end{align}
clear all beta coefficients and the rational constant in
\eqref{eq:forms}. The factors $2^{i+2}(i)_2$ only improve these
valuation bounds.

\subsection{Their cost after factorial normalization}
Legendre's formula gives
\begin{equation}\label{eq:legendre-lower}
 v_p(n!)\ge\frac{n}{p-1}-(\lambda_p+1).
\end{equation}
For each pole block $I_j(n)$, split the differences $b-a$ at zero.
Counting multiples of $p^k$ on the resulting two intervals gives
\[
 \sum_{\substack{b\in I_j(n)\\b\ne a}}v_p(b-a)
 \le\frac{\#I_j(n)}{p-1}+2(\lambda_p+1).
\]
After summing over the $31$ blocks,
\begin{equation}\label{eq:pole-cost}
 \sum_{b\ne a}B_bv_p(b-a)
 \le\frac{\gamma n+31}{p-1}+62(\lambda_p+1).
\end{equation}

For odd $p$, the integers $2a+1+2\ell$ consist of a negative
and a positive progression of odd integers, of lengths
\[
 A_-=(1+r)n-a+1,\qquad A_+=a+rn+1.
\]
The absolute values in either progression are $1,3,\ldots,2A-1$,
with $A=A_-$ or $A=A_+$. Their product is $(2A)!/(2^AA!)$, so
Legendre's formula gives
\[
 v_p\left(\frac{(2A)!}{2^AA!}\right)
       \ge\frac{A}{p-1}-(\lambda_p+1),
\]
since $2A<7n$. The lengths satisfy $A_-+A_+=Z/3$. Each product
occurs with multiplicity three, so the two estimates give
\[
 3\sum_{\ell\in\Lambda_n}v_p(2a+1+2\ell)
       \ge\frac{Z}{p-1}-6(\lambda_p+1).
\]
Combining this with \eqref{eq:pole-cost} yields
\[
 -\nu_p(a)\le\frac{Gn+25}{p-1}+68(\lambda_p+1).
\]
Equations \eqref{eq:small-ep} and \eqref{eq:legendre-lower} now give
\[
 e_p-Gv_p(n!)
 \le\frac{25}{p-1}+(68+33+G)(\lambda_p+1)
 \le128(\lambda_p+1).
\]
The maximum with zero in \eqref{eq:small-ep} causes no problem:
when it is attained at zero, the left side is nonpositive.
There are $O(\sqrt n)$ such primes, and
$(\lambda_p+1)\log p=O(\log n)$, so their total cost is $o(n)$.

At $2$, equation \eqref{eq:pole-cost} with $\lambda_2$ gives
\[
 -\nu_2(a)\le Z-1+\gamma n+31+62(\lambda_2+1)
            =(G+15)n+36+62(\lambda_2+1).
\]
Using $v_2(n!)\ge n-\lambda_2-1$ and \eqref{eq:two-ep}, we obtain
\[
 e_2-Gv_2(n!)\le15n+36+107(\lambda_2+1).
\]
Its contribution to $\log D_n-G\log(n!)$ is therefore at most
$15n\log2+o(n)$, as used in Section~\ref{sec:decay}.
